\section{Zbiory danych}
W badaniach wykorzystano zbiór JailbreakBench~\cite{chao2024jailbreakbenchopenrobustnessbenchmark} składający się szkodliwych zapytań. Zawiera 100 poleceń, które naruszają zasady bezpieczeństwa. Składa się z 55 promptów przygotowanych przez autorów zbioru, natomiast pozostałe 45 zostało dobrane z często stosowanych w literaturze zbiorów AdvBench~\cite{zou2023universaltransferableadversarialattacks} oraz HarmBench~\cite{mazeika2024harmbenchstandardizedevaluationframework}. Zapytania te podzielone są na 10 kategorii, które odzwierciedlają zasady użytkowania ustalone przez OpenAI~\cite{openai_usage_policy}.
Zbiór ten zawiera również 100 pytań bezpiecznych, które pozwoliły na ewaluację zjawiska nadmiernej odmowy modelu. Na ich podstawie sprawdzono również opóźnienia wprowadzone przez poszczególne metody obrony w stosunku do wariantu ich pozbawionego.
Ponadto, do wyznaczenia wektora odmowy użytego w metodzie obrony poprzez inżynierię aktywacji, wykorzystano zbiór pytań zamkniętych przygotowany przez Panickssery i in. w ramach badań nad CAA~\cite{panickssery2024steeringllama2contrastive}. Dodatkowo, do wyznaczenia w tym podejściu wektora sterującego bezpieczeństwem, wybrano 200 bezpiecznych i odpowiednio tyle samo zapytań szkodliwych ze zbioru XSTest~\cite{rottger2024xstesttestsuiteidentifying}. Ocenę użyteczności modeli oparto na teście MMLU~\cite{hendrycks2021measuringmassivemultitasklanguage} złożonego z pytań zamkniętych z zakresu 57 różnych dziedzin.